% ============================================================
\section{Positioning and Connections}
\label{sec:positioning}\lean{CoherenceCalculus.operatorCompression_eq_blockPP, CoherenceCalculus.unifiedGrammar_projection_algebra, CoherenceCalculus.towerCalculus_summary}
% ============================================================

This section situates Coherence Calculus relative to established mathematical
frameworks. Its ingredients are elementary: projection algebra, block
decomposition, and Hilbert-complex methods. It separates
standard background from the organizing contribution of the calculus.

\subsection{Finite observability in ergodic and systems settings}
\label{subsec:finite-observability-context}

The phrase \emph{finite observability} already has established narrow meanings
in nearby literatures. In ergodic theory, Ornstein and Weiss use it for
invariants that can be recovered almost surely from longer but still finite
typical observation windows, and show that entropy is the only finitely
observable invariant of stationary ergodic finite-valued processes in that
sense \cite{OrnsteinWeiss2007}. In linear systems theory, Kalman's
state-space framework and its descendants make observability precise by asking
which internal degrees of freedom are recoverable from finite output data
\cite{Kalman1963,Antoulas2005}. Coherence Calculus does not identify with
either framework, but it lives in the same conceptual neighborhood: the
observable/shadow split is treated as structural data rather than as a mere
numerical approximation device.

\subsection{Operator compression and projection-induced error}
\label{subsec:operator-compression}\lean{CoherenceCalculus.operatorCompression_eq_blockPP, CoherenceCalculus.leakageCorrectionTerm}

The basic operation of Coherence Calculus, namely compressing an operator $A$
to $\CCRestr{A}{h} = \CCPi{h} A \CCPi{h}$, is standard in numerical linear
algebra and model reduction.
Projection-based compression and the resulting coarse/fine correction terms are
standard in reduced-order modeling, multilevel solvers, and open-system
effective dynamics \cite{Antoulas2005,BriggsHensonMcCormick2000,BreuerPetruccione2002}.
The ``leakage correction'' $\CCLeakCorr{h}{A}{B} = \CCPi{h} A \CCPibar{h} B \CCPi{h}$ that appears in product compression is the same cross-term that arises in:
\begin{itemize}[itemsep=0.2em]
\item \textbf{Reduced-order modeling:} Galerkin projection onto a subspace introduces closure terms from the neglected modes.
\item \textbf{Multigrid and domain decomposition:} Coarse-grid corrections require accounting for fine-grid residuals.
\item \textbf{Quantum open systems:} Tracing over environment degrees of freedom generates effective dynamics with ``system-bath'' correction terms.
\end{itemize}
Coherence Calculus organizes these correction terms into a common grammar of
chain rules, product rules, and elimination patterns, rather than treating
each application separately.

\subsection{Conditional expectation and coarse-graining}
\label{subsec:coarse-graining}\lean{CoherenceCalculus.shadow_plus_proj_identity, CoherenceCalculus.proj_shadow_orthogonal}

When $\CCPi{h}$ is a conditional expectation (e.g., block averaging), the coherence defect $\CCTau{h}{F}{x} = \CCPi{h} F(x) - F(\CCPi{h} x)$ is the ``closure term'' familiar from:
\begin{itemize}[itemsep=0.2em]
\item \textbf{Turbulence modeling:} Reynolds decomposition and subgrid-scale closures.
\item \textbf{Statistical mechanics:} Coarse-grained free energies and renormalization group.
\item \textbf{Large eddy simulation:} Filtered equations require modeling of unresolved scales.
\end{itemize}
These interpretations are standard in coarse-grained fluid and statistical
descriptions \cite{Pope2000,Goldenfeld1992}.
The bilinear defect $\CCTauBi{h}{B}{x}{y}$
(Proposition~\ref{prop:product-rule}) is the same algebraic object that
appears as subgrid stress or eddy-viscosity defect in these settings.
Coherence Calculus does not solve the closure problem; it records exact
identities for the defect structure. The broader idea that elimination of
unresolved variables produces memory, noise, and effective closure terms is
also classical in the Mori--Zwanzig and optimal-prediction literature
\cite{Zwanzig1961,Mori1965,ChorinHaldKupferman2000}.

\subsection{Schur complement and Feshbach projection}
\label{subsec:schur-feshbach}\lean{schurComplement, effectiveOp, schur_factorization}

The Schur complement formula (Section~\ref{sec:resolvent-defect}) is classical
\cite{Haynsworth1968,Zhang2005}:
\[
(M/D) = A - BD^{-1}C \quad \text{for } M = \begin{pmatrix} A & B \\ C & D \end{pmatrix}.
\]
In physics, this is known as the \emph{Feshbach projection} or \emph{effective Hamiltonian} method
\cite{Feshbach1958,Feshbach1962}:
\begin{itemize}[itemsep=0.2em]
\item \textbf{Quantum mechanics:} Projecting onto a ``P-space'' generates an effective operator with ``self-energy'' corrections from the ``Q-space.''
\item \textbf{Scattering theory:} T-matrix formulations via elimination of closed channels.
\item \textbf{Renormalization:} Integrating out high-energy modes generates effective low-energy couplings.
\item \textbf{Statistical coarse-graining:} Projection formalisms generate exact memory and return terms for unresolved variables \cite{Zwanzig1961,Mori1965,ChorinHaldKupferman2000}.
\end{itemize}
Coherence Calculus uses the Schur complement as the canonical elimination
pattern for shadow degrees of freedom. The distinctive step is its systematic
connection to the tower calculus through derivative, antiderivative, and
telescoping identities.

\subsection{Multiresolution and nested projections}
\label{subsec:multiresolution}\lean{CoherenceCalculus.multiresolution_nesting, CoherenceCalculus.wavelet_detail_analog, CoherenceCalculus.increment_orthogonality}

The horizon tower $\CCPi{0} \le \CCPi{1} \le \cdots$ is a nested sequence of orthogonal projections.
This structure appears in:
\begin{itemize}[itemsep=0.2em]
\item \textbf{Wavelet analysis:} Multiresolution analysis with nested approximation spaces $V_0 \subset V_1 \subset \cdots$.
\item \textbf{Martingale theory:} Filtrations and conditional expectations with respect to nested $\sigma$-algebras.
\item \textbf{Spectral theory:} Spectral projections $E_{(-\infty, \lambda]}$ for self-adjoint operators.
\end{itemize}
Standard references include multiresolution analysis, martingale filtrations,
and spectral projection theory \cite{Mallat1989,Daubechies1992,Doob1953,ReedSimon1972}.
The tower derivative $\CCDer{A} = [\CCIndexOp, A]$ and antiderivative
$\CCInt{B}$ provide algebraic tools for these nested structures. The
fundamental theorem (Theorem~\ref{thm:ftc-tower}) is an operator analogue of a
reconstruction formula.

\subsection{Hilbert complexes and Hodge theory}
\label{subsec:hodge}\lean{CoherenceCalculus.complex_nilpotency, CoherenceCalculus.intertwining_preserves_complex, CoherenceCalculus.cohomology_splits}

The solver layer (Section~\ref{subsec:solver-layer}) uses standard finite-dimensional Hodge theory
\cite{ArnoldFalkWinther2006,Schwarz1995,Warner1983}:
\begin{itemize}[itemsep=0.2em]
\item \textbf{Hodge decomposition:} $H^k = \Img(D_{k-1}) \oplus \Img(D_k^*) \oplus \Ker(\mathcal{L}_k)$ is the standard orthogonal decomposition.
\item \textbf{Green operator:} $\CCGreen{k} = \mathcal{L}_k^{-1}$ under exactness is the standard pseudo-inverse.
\item \textbf{Homotopy identity:} $DK + KD = I$ under exactness is the standard Hodge--de Rham result.
\end{itemize}
For infinite-dimensional Hilbert complexes, the same circle of ideas extends to
functional-analytic Hilbert-complex settings and finite element exterior
calculus \cite{ArnoldFalkWinther2006}.
Coherence Calculus restricts to the finite-dimensional case at the structural
level. Infinite-dimensional extensions require the usual functional-analytic
care with closedness and Fredholm conditions.

\subsection{Distinctive Features}
\label{subsec:what-is-new}\lean{CoherenceCalculus.unifiedGrammar_projection_algebra, CoherenceCalculus.towerCalculus_summary, CoherenceCalculus.diagonal_offdiagonal_decomposition}

The individual components above are standard. In the present setting, the
distinctive features are:
\begin{enumerate}[label=(\roman*), itemsep=0.3em]
\item \textbf{Unified grammar:} Chain rule, product rule, and elimination patterns organized around the defect identity.
\item \textbf{Tower calculus:} Derivative $\partial = [\CCIndexOp, \cdot]$, antiderivative $\CCInt{\cdot}$, fundamental theorem, and IBP as a self-contained calculus on operator algebras.
\item \textbf{Completion-and-selection algebra:} Systematic reduction of locality, symmetry, and conservation constraints to explicit operator selections via isotypic decomposition.
\item \textbf{Integration of solver layer:} Hodge/Green/homotopy operators connected to the tower structure, with holographic saturation as relative exactness.
\end{enumerate}
